% GENERATED FILE. Edit canonical lessons, then run npm run generate:books.
% canonical-source-hash: fnv1a64:6f553fcab1866183
% canonical-lessons: MR-C435-tancai, MR-C435-tutvada, MR-C435-korda, MR-C435-atne, MR-C435-nasne, MR-R435-review-428-435

\chapter{Shortages}
\label{ch:mr-shortages}
\hlchaptermodality{435}

\begin{chapteropening}
\textbf{By the end of this chapter:} \emph{I can talk about shortages and scarcity, and say a well dried up or food went bad.}

Complete the last lesson of chapter 435: I can talk about shortages and scarcity, and say a well dried up or food went bad.
\end{chapteropening}

\section[\texorpdfstring{\mr{टंचाई} (ṭañcāī)}{टंचाई (ṭañcāī)}]{\mr{टंचाई} (ṭañcāī) — a shortage}

\label{lesson:MR-C435-tancai}



\begin{quote}
Before the new one: say the Marathi for to rush about, then the Marathi for to find a way.
\end{quote}

\subsection*{You'll want to know: \mr{टंचाई}}
\textbf{\mr{टंचाई}} — \emph{ṭañcāī} — \textquotedblleft{}a shortage\textquotedblright{}.

\begin{quote}
\mr{उन्हाळ्यात} \mr{गावात} \mr{पाण्याची} \mr{टंचाई} \mr{असते}.
\end{quote}

\begin{quote}
\emph{unhāḷyāt gāvāt pāṇyācī ṭañcāī aste.}
\end{quote}

\begin{quote}
In summer the village is short of water.
\end{quote}

Feminine.

The thing in short supply takes -\mr{ची}: \mr{पाण्याची} \mr{टंचाई}.

\begin{cousinweb}
\textbf{\mr{टंचाई}} (\emph{ṭañcāī}) is a Marathi word for a shortage of water, food or money; its origin is not settled.
\end{cousinweb}

\begin{grammarlens}[title={What you've built}]
One more word for getting around.
\end{grammarlens}

\subsection*{Your turn}
\begin{itemize}
\raggedright
  \item \emph{Say it:} \emph{ṭañcāī}
  \item \emph{Say it:} \emph{ṭañcāī}, once more
  \item \emph{Say it:} \emph{unhāḷyāt gāvāt pāṇyācī ṭañcāī aste.} — In summer the village is short of water.
  \item \emph{Recall:} say the Marathi for to rush about, then the Marathi for to find a way, then say \emph{ṭañcāī} again
\end{itemize}

\begin{tcolorbox}[breakable,colback=teal!4,colframe=teal!35!black,title={Before you move on}]
What does \mr{टंचाई} mean? (\textquotedblleft{}A shortage\textquotedblright{}.) Say it once more.
\end{tcolorbox}
\section[\texorpdfstring{\mr{तुटवडा} (tuṭvaḍā)}{तुटवडा (tuṭvaḍā)}]{\mr{तुटवडा} (tuṭvaḍā) — a scarcity}

\label{lesson:MR-C435-tutvada}



\begin{quote}
Before the new one: say the Marathi for to find a way, then the Marathi for a shortage.
\end{quote}

\subsection*{You'll want to know: \mr{तुटवडा}}
\textbf{\mr{तुटवडा}} — \emph{tuṭvaḍā} — \textquotedblleft{}a scarcity\textquotedblright{}.

\begin{quote}
\mr{बाजारात} \mr{कांद्याचा} \mr{तुटवडा} \mr{आहे}.
\end{quote}

\begin{quote}
\emph{bājārāt kāndyācā tuṭvaḍā āhe.}
\end{quote}

\begin{quote}
There is a scarcity of onions in the market.
\end{quote}

Masculine.

Built on \mr{तुट}-, the stem of \mr{तुटणे}.

\begin{cousinweb}
\textbf{\mr{तुटवडा}} (\emph{tuṭvaḍā}) is built on the stem of \textbf{\mr{तुटणे}}, \textquotedblleft{}to break, to fall short\textquotedblright{} (CDIAL 6065), with a noun ending.
\end{cousinweb}

\begin{grammarlens}[title={What you've built}]
One more word for getting around.
\end{grammarlens}

\subsection*{Your turn}
\begin{itemize}
\raggedright
  \item \emph{Say it:} \emph{tuṭvaḍā}
  \item \emph{Say it:} \emph{tuṭvaḍā}, once more
  \item \emph{Say it:} \emph{bājārāt kāndyācā tuṭvaḍā āhe.} — There is a scarcity of onions in the market.
  \item \emph{Recall:} say the Marathi for to find a way, then the Marathi for a shortage, then say \emph{tuṭvaḍā} again
\end{itemize}

\begin{tcolorbox}[breakable,colback=teal!4,colframe=teal!35!black,title={Before you move on}]
What does \mr{तुटवडा} mean? (\textquotedblleft{}A scarcity\textquotedblright{}.) Say it once more.
\end{tcolorbox}
\section[\texorpdfstring{\mr{कोरडा} (korḍā)}{कोरडा (korḍā)}]{\mr{कोरडा} (korḍā) — parched; dry}

\label{lesson:MR-C435-korda}



\begin{quote}
Before the new one: say the Marathi for a shortage, then the Marathi for a scarcity.
\end{quote}

\subsection*{You'll want to know: \mr{कोरडा}}
\textbf{\mr{कोरडा}} — \emph{korḍā} — \textquotedblleft{}parched; dry\textquotedblright{}.

\begin{quote}
\mr{या} \mr{वर्षी} \mr{विहीर} \mr{कोरडी} \mr{पडली}.
\end{quote}

\begin{quote}
\emph{yā varṣī vihīr korḍī paḍlī.}
\end{quote}

\begin{quote}
This year the well went dry.
\end{quote}

\mr{कोरडा} \mr{पडणे}, to go dry; \mr{कोरडी}, \mr{कोरडं} for the other genders.

\begin{cousinweb}
\textbf{\mr{कोरडा}} (\emph{korḍā}) is a Marathi adjective for land, wells or weather without water; its origin is not settled.
\end{cousinweb}

\begin{grammarlens}[title={What you've built}]
One more word for getting around.
\end{grammarlens}

\subsection*{Your turn}
\begin{itemize}
\raggedright
  \item \emph{Say it:} \emph{korḍā}
  \item \emph{Say it:} \emph{korḍā}, once more
  \item \emph{Say it:} \emph{yā varṣī vihīr korḍī paḍlī.} — This year the well went dry.
  \item \emph{Recall:} say the Marathi for a shortage, then the Marathi for a scarcity, then say \emph{korḍā} again
\end{itemize}

\begin{tcolorbox}[breakable,colback=teal!4,colframe=teal!35!black,title={Before you move on}]
What does \mr{कोरडा} mean? (\textquotedblleft{}Parched; dry\textquotedblright{}.) Say it once more.
\end{tcolorbox}
\section[\texorpdfstring{\mr{आटणे} (āṭṇe)}{आटणे (āṭṇe)}]{\mr{आटणे} (āṭṇe) — to dry up}

\label{lesson:MR-C435-atne}



\begin{quote}
Before the new one: say the Marathi for a scarcity, then the Marathi for parched.
\end{quote}

\subsection*{You'll want to know: \mr{आटणे}}
\textbf{\mr{आटणे}} — \emph{āṭṇe} — \textquotedblleft{}to dry up\textquotedblright{}.

\begin{quote}
\mr{उन्हाळ्यात} \mr{नदी} \mr{आटली}.
\end{quote}

\begin{quote}
\emph{unhāḷyāt nadī āṭlī.}
\end{quote}

\begin{quote}
In summer the river dried up.
\end{quote}

Past (he …): \mr{तो} \mr{आटला}, \emph{to āṭlā}.

For water in a pot too: \mr{दूध} \mr{आटलं}, the milk boiled down.

\begin{cousinweb}
\textbf{\mr{आटणे}} (\emph{āṭṇe}) is a Marathi verb for water drying away; its origin is not settled.
\end{cousinweb}

\begin{grammarlens}[title={What you've built}]
One more word for getting around.
\end{grammarlens}

\subsection*{Your turn}
\begin{itemize}
\raggedright
  \item \emph{Say it:} \emph{āṭṇe}
  \item \emph{Say it:} \emph{āṭṇe}, once more
  \item \emph{Say it:} \emph{unhāḷyāt nadī āṭlī.} — In summer the river dried up.
  \item \emph{Recall:} say the Marathi for a scarcity, then the Marathi for parched, then say \emph{āṭṇe} again
\end{itemize}

\begin{tcolorbox}[breakable,colback=teal!4,colframe=teal!35!black,title={Before you move on}]
What does \mr{आटणे} mean? (\textquotedblleft{}To dry up\textquotedblright{}.) Say it once more.
\end{tcolorbox}
\section[\texorpdfstring{\mr{नासणे} (nāsṇe)}{नासणे (nāsṇe)}]{\mr{नासणे} (nāsṇe) — to go bad}

\label{lesson:MR-C435-nasne}



\begin{quote}
Before the new one: say the Marathi for parched, then the Marathi for to dry up.
\end{quote}

\subsection*{You'll want to know: \mr{नासणे}}
\textbf{\mr{नासणे}} — \emph{nāsṇe} — \textquotedblleft{}to go bad\textquotedblright{}.

\begin{quote}
\mr{उन्हामुळे} \mr{दूध} \mr{नासलं}.
\end{quote}

\begin{quote}
\emph{unhāmuḷe dūdh nāslaṁ.}
\end{quote}

\begin{quote}
The milk went bad in the heat.
\end{quote}

Past (he …): \mr{तो} \mr{नासला}, \emph{to nāslā}.

Milk, fruit and cooked food \mr{नासतात}.

\begin{cousinweb}
\textbf{\mr{नासणे}} (\emph{nāsṇe}) is inherited: CDIAL lists Marathi \emph{nāsṇẽ}, \textquotedblleft{}to be spoilt\textquotedblright{}, under Sanskrit \emph{naśyati}, \textquotedblleft{}is lost, perishes\textquotedblright{} (CDIAL 7027).
\end{cousinweb}

\begin{grammarlens}[title={What you've built}]
One more word for getting around.
\end{grammarlens}

\subsection*{Your turn}
\begin{itemize}
\raggedright
  \item \emph{Say it:} \emph{nāsṇe}
  \item \emph{Say it:} \emph{nāsṇe}, once more
  \item \emph{Say it:} \emph{unhāmuḷe dūdh nāslaṁ.} — The milk went bad in the heat.
  \item \emph{Recall:} say the Marathi for parched, then the Marathi for to dry up, then say \emph{nāsṇe} again
\end{itemize}

\begin{tcolorbox}[breakable,colback=teal!4,colframe=teal!35!black,title={Before you move on}]
What does \mr{नासणे} mean? (\textquotedblleft{}To go bad\textquotedblright{}.) Say it once more.
\end{tcolorbox}
\section[(dialogue)]{Looking back: from Spending and Saving to Shortages}

\label{lesson:MR-R435-review-428-435}



\begin{quote}
Say the Marathi for to squander and a balance.
\end{quote}

\subsection*{Your turn}
Say these aloud:

\begin{itemize}
\raggedright
  \item to squander, a balance, an investment, what one can afford, economical
  \item an auction, collateral, to pawn, to recover money, a lease
  \item a payment in return, a donation, a subscription, a share, to be left over
  \item a breakdown, out of action, to snap, to burst, to rip
  \item a hole, a small hole, a stain, to get messed up, to ooze
  \item a scrape, a split, to rust, to get dented, to stop working
  \item inconvenience, missing each other, a fiasco, to rush about, to find a way
  \item a shortage, a scarcity, parched, to dry up, to go bad
\end{itemize}

\begin{tcolorbox}[breakable,colback=teal!4,colframe=teal!35!black,title={Before you move on}]
To squander? (\textbf{\mr{उधळणे}}, \emph{udhaḷṇe}.) A balance? (\textbf{\mr{शिल्लक}}, \emph{śillak}.) An investment? (\textbf{\mr{गुंतवणूक}}, \emph{guntavṇūk}.) What one can afford? (\textbf{\mr{ऐपत}}, \emph{aipat}.) Economical? (\textbf{\mr{किफायतशीर}}, \emph{kiphāyatśīr}.) An auction? (\textbf{\mr{लिलाव}}, \emph{lilāv}.) Collateral? (\textbf{\mr{तारण}}, \emph{tāraṇ}.) To pawn? (\textbf{\mr{गहाण} \mr{ठेवणे}}, \emph{gahāṇ ṭhevṇe}.) To recover money? (\textbf{\mr{वसूल} \mr{करणे}}, \emph{vasūl karṇe}.) A lease? (\textbf{\mr{भाडेपट्टा}}, \emph{bhāḍepaṭṭā}.) A payment in return? (\textbf{\mr{मोबदला}}, \emph{mobadlā}.) A donation? (\textbf{\mr{देणगी}}, \emph{deṇgī}.) A subscription? (\textbf{\mr{वर्गणी}}, \emph{vargaṇī}.) A share? (\textbf{\mr{हिस्सा}}, \emph{hissā}.) To be left over? (\textbf{\mr{उरणे}}, \emph{urṇe}.) A breakdown? (\textbf{\mr{बिघाड}}, \emph{bighāḍ}.) Out of action? (\textbf{\mr{निकामी}}, \emph{nikāmī}.) To snap? (\textbf{\mr{तुटणे}}, \emph{tuṭṇe}.) To burst? (\textbf{\mr{फुटणे}}, \emph{phuṭṇe}.) To rip? (\textbf{\mr{फाटणे}}, \emph{phāṭṇe}.) A hole? (\textbf{\mr{भोक}}, \emph{bhok}.) A small hole? (\textbf{\mr{छिद्र}}, \emph{chidra}.) A stain? (\textbf{\mr{डाग}}, \emph{ḍāg}.) To get messed up? (\textbf{\mr{विस्कटणे}}, \emph{viskaṭṇe}.) To ooze? (\textbf{\mr{पाझरणे}}, \emph{pājharṇe}.) A scrape? (\textbf{\mr{ओरखडा}}, \emph{orakhḍā}.) A split? (\textbf{\mr{तडा}}, \emph{taḍā}.) To rust? (\textbf{\mr{गंजणे}}, \emph{gañjṇe}.) To get dented? (\textbf{\mr{चेपणे}}, \emph{cepṇe}.) To stop working? (\textbf{\mr{बंद} \mr{पडणे}}, \emph{band paḍṇe}.) Inconvenience? (\textbf{\mr{गैरसोय}}, \emph{gairsoy}.) Missing each other? (\textbf{\mr{चुकामूक}}, \emph{cukāmūk}.) A fiasco? (\textbf{\mr{फजिती}}, \emph{phajitī}.) To rush about? (\textbf{\mr{धावपळ} \mr{करणे}}, \emph{dhāvpaḷ karṇe}.) To find a way? (\textbf{\mr{मार्ग} \mr{काढणे}}, \emph{mārga kāḍhṇe}.) A shortage? (\textbf{\mr{टंचाई}}, \emph{ṭañcāī}.) A scarcity? (\textbf{\mr{तुटवडा}}, \emph{tuṭvaḍā}.) Parched? (\textbf{\mr{कोरडा}}, \emph{korḍā}.) To dry up? (\textbf{\mr{आटणे}}, \emph{āṭṇe}.) To go bad? (\textbf{\mr{नासणे}}, \emph{nāsṇe}.)
\end{tcolorbox}
